\documentclass[10pt,a4paper]{amsart}
\usepackage[margin=19mm]{geometry}
\usepackage{amsmath,amssymb,mathtools}
\usepackage[scr=rsfs,cal=euler]{mathalfa}
\usepackage{xfrac}
\usepackage[unicode,hidelinks,bookmarks=false]{hyperref}
\newcommand{\ie}{i.e.\ }
\newcommand{\cf}{cf.\ }
\newcommand{\resp}{respectively\ }
\newcommand{\Subsection}{\subsection}
\providecommand{\nobreakdash}{\nobreak\hbox{-}}
\setlength{\emergencystretch}{2em}
\multlinegap0pt
\usepackage{bm}
\usepackage{bbm}
\usepackage{dsfont}
\usepackage{xspace}
\usepackage{cancel}
\usepackage{mathtools}
\usepackage{amsthm}
\makeatletter
\@ifundefined{lem}{\newtheorem{lem}{Lemma}}{}
\@ifundefined{prop}{\newtheorem{prop}[lem]{Proposition}}{}
\makeatother
\DeclareMathOperator{\im}{Im}
\DeclareMathOperator{\re}{Re}
\newcommand{\ph}{\varphi}
\newcommand{\z}{\zeta}
\title[The Minimal Absolute Value of Sums of Fifth Roots of Unity]{The Minimal Absolute Value of Sums of Fifth Roots of Unity}
\author{Akihiro Munemasa, Guillermo N\'u\~nez Ponasso}
\date{Source: arXiv:2607.00825v2 (CC BY 4.0)}
\begin{document}
\maketitle

\noindent\emph{Source and licence.} The Minimal Absolute Value of Sums of Fifth Roots of Unity. Version arXiv:2607.00825v2 (CC BY 4.0), distributed under the Creative Commons Attribution 4.0 International licence, as recorded in the arXiv licence metadata for this exact version and shown on the abstract page https://arxiv.org/abs/2607.00825v2. Full rights notice: licenses/arxiv-2607.00825-CC-BY-4.0.txt.\par\smallskip

\noindent\textit{Editorial scope.} Counted unit: the Prop whose source label is \texttt{prop:1f} in the article (source0.tex lines 820--825, source label \texttt{prop:1f}), delivered together with the article's own complete original proof of it (source0.tex lines 826--874, one \texttt{proof} environment of 49 lines). No mathematics is written, repaired or re-derived: statement and proof are the article's own text line for line. Required context retained uncounted: lem:2 (lines 398--404); lem:bnxy (lines 425--432); lem:ub2 (lines 648--650). Every definition, display and label of those lines is retained unchanged. Omitted from this package and not redistributed: 1--397, 405--424, 433--647, 651--819, 875--899. Nothing inside a retained excerpt is omitted or altered: every retained body is delivered exactly as the article's lines are written, so no omission receipt of any kind accompanies this package. Citations: the retained excerpts carry no citation command at all, so no bibliography entry of the article is transcribed and no reference is redistributed. Reference adaptation: none was needed and none was made; every reference inside the delivered unit resolves inside it, so no unresolved reference or citation is produced. Printed slips kept literal: no printed word, symbol, hypothesis or number of the counted statement or of its proof was altered. Layout and class: the article's own class is \texttt{article} with packages including geometry, graphicx, fontenc, tgpagella, footmisc, hyperref, amsmath, amssymb, amsthm, rotating, multirow, booktabs, adjustbox, mathtools; the wrapper is the lane's pinned \texttt{amsart} editorial wrapper at 10pt on A4, which restates the theorem-like environments the retained text needs and adds the article's own macro definitions that the retained text needs (\texttt{\textbackslash im}, \texttt{\textbackslash ph}, \texttt{\textbackslash re}, \texttt{\textbackslash z}), with extra wrapper packages (bm, bbm, dsfont, xspace, cancel, mathtools). The delivered numbering is the unit's own. Assets: no retained span contains a figure environment, a graphics-inclusion command, a PDF-page inclusion, a table or a TikZ picture; no image file is copied into this package, the package declares no asset dependency and no image file is redistributed with it. Licence: version arXiv:2607.00825v2 of this article is distributed under the Creative Commons Attribution 4.0 International licence, as recorded in the arXiv licence metadata for this exact version and shown on the abstract page https://arxiv.org/abs/2607.00825v2. A full rights notice is delivered at licenses/arxiv-2607.00825-CC-BY-4.0.txt. Characters: every retained excerpt is pure ASCII, so the delivered engine, XeLaTeX with its default Latin Modern text font, reports no missing character.
\begin{lem}\label{lem:2} Let $t=\ph$ or $\alpha$. If $b>0$ is an integer with Ostrowski representation
  $b=\sum_{k=r}^{R}c_{k}F_{k}$,
 then
  \begin{equation}\label{eq:lem2}
    \|bt\|>\frac{1}{\ph^{r+1}}.
  \end{equation}
\end{lem}

\begin{lem}\label{lem:bnxy}
    Suppose that  \[s=a_1\z+a_2\z^2+a_3\z^3+a_4\z^4\in\Sigma_5^0(n).\]
    Then,
    \[|\re(s)|=\frac{1}{2\ph}|n-b\alpha|,\]
    where $b:=a_2+a_3$, and
    \[|\im(s)|^2=\frac{4-\ph^2}{4}|x\ph- y|^2,\]
    where $x:=a_1-a_4$, and $y:=a_3-a_2$. Furthermore,  $1/\ph^{4}\geq |s|$ implies that $n=[b\alpha]$ and $b=[n/\alpha]$; and $1/\ph^3\geq |s|$ implies that $y=[x\ph]$ and $x=[y/\ph]$.
\end{lem}

\begin{lem} \label{lem:ub2} Let $k\geq 1$ be an integer and let $n=5k$. Then
  \[\sigma_5(n)^2\leq \frac{\sqrt{5}}{\ph^{2\nu(k)-1}}.\]
\end{lem}

\begin{prop}\label{prop:1f} Let $m\geq 5$ and $n=5F_{m}$. If \[s=a_1\z+a_2\z^2+a_3\z^3+a_4\z^4\in\Sigma_5^0(n),\]
 and $a_2+a_3=L_{m-1}$, then
  \[|s|^2\geq \frac{\sqrt{5}}{\ph^{2m-1}}.\]
 Furthermore, $s$ satisfies the above bound with equality if and only if $s$ or $\overline{s}$ is equal to
  \[2F_m\z + F_{m-2}\z^2 +F_m\z^3+F_{m+1}\z^4.\]
\end{prop}

\begin{proof}
We can assume
  \[|s|^2\leq \frac{\sqrt{5}}{\ph^{2m-1}}\leq \frac{\sqrt{5}}{\ph^9}<\frac{1}{\ph^4}.\]
  Let $b:=a_2+a_3$. Applying Lemma~\ref{lem:bnxy}, and using the fact that $n=L_{m+1}+L_{m-1}$ and $b=L_{m-1}=F_{m}+F_{m-2}$, we have
  \begin{align*}
    |\re(s)|^2 &= \frac{1}{4\ph^2}|n-b\alpha|^2=\frac{5}{4\ph^{2m}},
  \end{align*}
  and
  \[|\im(s)|^2=\frac{4-\ph^2}{4}|x\ph-y|^2,\]
  where $x=a_1-a_4$, and $y=a_3-a_2$ satisfy $y=[x\ph]$ and $x=[y/\ph]$. Noting that $a_1+a_4=5F_{m}-b=L_{m+1}$, we can express the coefficients of $s$ in terms of $x$ and $y$ as
  \[(a_1,a_2,a_3,a_4)=\frac{1}{2}(L_{m+1}+x,L_{m-1}-y,L_{m-1}+y,L_{m+1}-x).\]
   Also, from $L_{m+1}\not\equiv L_{m-1}\pmod{2}$ we find that $x=0$ implies $y\neq 0$. If $x=0$, then for $m\geq 3$,
  \begin{align*}
    |s|^2\geq |\im(s)|^2\geq\frac{4-\ph^2}{4}>\frac{\sqrt{5}}{\ph^{2m-1}},
  \end{align*}
 Hence, without loss of generality, we may assume that $x>0$. Observe that
 \[\frac{\sqrt{5}}{\ph^{2m-1}}\geq |s|^2\geq \frac{5}{4\ph^{2m}}+\frac{4-\ph^2}{4}\|x\ph\|^2.\]
 So, it follows that $\|x\ph\|^2\leq 1/\ph^{2m-4}$. Let $x=\sum_{k=r}^R c_{k}F_{k}$ be the Ostrowski representation of $x$. Then, Lemma~\ref{lem:2} implies $\|x\ph\|^2>1/\ph^{2r+2}$, so $r\geq m-2$ and, in particular, $x\geq F_{m-2}$. From $0\leq 2a_2=L_{m-1}-y$, it follows that $y\leq L_{m-1}=F_{m}+F_{m-2}$, and then, for $m\geq 5$, $x=[y/\ph]\leq F_{m-1}+F_{m-3}$. Hence, the only valid Ostrowski representations of $x$ are $x= F_{m-2}$, or $x=F_{m-1}$. 
%  \begin{equation*}
%  \begin{tabular}{rllll}
%    $x=$ & $F_{m-2}$, & $F_{m}$, &  $F_{m+1}$,& $F_{m+2}$,\\
         %& $F_{m-1}$, & $F_{m}+F_{m-3}$, & $F_{m+1}+F_{m-3}$,& $F_{m+2}+F_{m-3}$,\\
         %& $F_{m-1}$, & $F_{m}+F_{m-2}$ & $F_{m+1}+F_{m-2}$, & $F_{m+2}+F_{m-2}$,\\
         %& & & $F_{m+1}+F_{m-1}$, & $F_{m+2}+F_{m-1}$,\\
         %& & & $F_{m+1}+F_{m-1}+F_{m-3}$,& $F_{m+2}+F_{m-1}+F_{m-3}$, \\
         %& & & & $F_{m+2}+F_{m}$.
%  \end{tabular}
%  \end{equation*}
 % Since $m\geq 5$ we have $r\geq m-2\geq 3$, so Lemma~\ref{lem:1} gives $y=[x\ph]=\sum_{k=r}^{R}c_{k}F_{k+1}$. If  $x\geq F_{m}$, then $y\geq F_{m+1}$, and $2a_2\leq L_{m-1}-F_{m+1}=-F_{m-3}<0$, a contradiction. Hence $x\in\{F_{m-2},F_{m-1}\}$.
  
   Case 1:  $x=F_{m-1}$. Then  $y=F_{m}$, $2a_2=L_{m-1}-F_{m}=F_{m-2}$, and $2a_4=L_{m+1}-F_{m-1}=3F_{m}$. But this implies $F_{m-2}\equiv F_{m}\equiv 0\pmod{2}$, a contradiction.% So this case admits no integral solution for $(a_1,a_2,a_3,a_4)$.
  
  %\[x=a_1-a_4\equiv a_1+a_4=L_{m+1}\equiv F_{m+1}\pmod{2}.\]
  %The only values of $x$ in the list above that satisfy $x\equiv F_{m+1}\pmod{2}$ are \[x=F_{m-2},F_{m+2}+F_{m}.%F_{m-1}+F_{m-3},F_{m+1},F_{m+2}+F_{m-3}
  %\]
  %However, the system
  %\begin{equation*}\label{eq:sys1}
  %\begin{cases}
  %  a_1+a_4 = L_{m+1}\\
  %  a_1-a_4 = x\\
  %  a_2+a_3 = L_{m-1}\\
  %  a_3-a_2 = [x\ph]
  %\end{cases}
  %\end{equation*}
  %has a solution $(a_1,a_2,a_3,a_4)\in\N^4$ if and only 
  Case 2: $x=F_{m-2}$. Then $y=F_{m-1}$, and  $(a_1,a_2,a_3,a_4)=(2F_{m},F_{m-2},F_m,F_{m+1})$. The corresponding sum $s$ satisfies $|s|^2=\sqrt{5}/\ph^{2m-1}$, as shown in the proof of Lemma~\ref{lem:ub2}.
  %
  %If $x=F_{m-1}+F_{m-3}=L_{m-2}$ then $(a_1,a_2,a_3,a_4)=(L_{m},0,L_{m-1},L_{m-1})$,  and in this case  $|s|^2 =\frac{5}{\ph^{2m-2}}>\frac{\sqrt{5}}{\ph^{2m-1}}$.
\end{proof}

\end{document}
